\documentclass[
    language=english,
    doctype=article,
    institution=none,
]{../../omnilatex}

\RequirePackage{config/document-settings}
\addbibresource{../../bib/bibliography.bib}

\begin{document}

\maketitle

\section{Scope and Purpose}

This regulation establishes minimum requirements for the collection,
processing, storage, and transmission of personal data by all entities
operating within the jurisdiction of the National Data Protection Authority
(NDPA). It applies to:

\begin{itemize}
    \item All public-sector organizations that process personal data.
    \item Private-sector organizations that process personal data of more
        than 10,000 individuals annually.
    \item Data processors acting on behalf of the above organizations.
    \item Cross-border data transfers originating from within the
        jurisdiction.
\end{itemize}

The purpose of this regulation is to protect the fundamental rights and
freedoms of natural persons with regard to the processing of their personal
data, while ensuring the free flow of data within the jurisdiction.

\section{Definitions}

\begin{description}
    \item[``Personal Data''] means any information relating to an identified
        or identifiable natural person, including but not limited to names,
        identification numbers, location data, online identifiers, and
        factors specific to the physical, physiological, genetic, mental,
        economic, cultural, or social identity of that person.
    \item[``Data Controller''] means the natural or legal person which
        determines the purposes and means of the processing of personal
        data.
    \item[``Data Processor''] means a natural or legal person which
        processes personal data on behalf of the controller.
    \item[``Processing''] means any operation performed on personal data,
        including collection, recording, organization, storage, adaptation,
        retrieval, consultation, use, disclosure, dissemination, alignment,
        combination, restriction, erasure, or destruction.
    \item[``Data Breach''] means a breach of security leading to the
        accidental or unlawful destruction, loss, alteration, unauthorized
        disclosure of, or access to, personal data.
    \item[``Data Protection Officer'' (DPO)] means the individual
        designated by an organization to oversee compliance with this
        regulation.
\end{description}

\section{Data Processing Requirements}

\subsection{Lawful Basis for Processing}

Organizations shall process personal data only where one or more of the
following lawful bases apply:

\begin{enumerate}
    \item The data subject has given consent to the processing.
    \item Processing is necessary for the performance of a contract.
    \item Processing is necessary for compliance with a legal obligation.
    \item Processing is necessary to protect the vital interests of the
        data subject.
    \item Processing is necessary for the performance of a task carried out
        in the public interest.
    \item Processing is necessary for the legitimate interests pursued by
        the controller, except where overridden by the interests, rights,
        and freedoms of the data subject.
\end{enumerate}

\subsection{Data Minimization}

Organizations shall ensure that personal data collected is adequate,
relevant, and limited to what is necessary in relation to the purposes for
which it is processed.

\subsection{Storage Limitation}

Personal data shall be kept in a form that permits identification of data
subjects for no longer than is necessary for the purposes for which the
data is processed. Organizations must document and justify all retention
periods.

\section{Compliance Requirements}

The following table summarizes the key compliance obligations and their
deadlines:

\begin{table}[h]
    \centering
    \begin{tabular}{lllr}
        \toprule
        \textbf{Requirement} & \textbf{Obligation} & \textbf{Deadline} & \textbf{Penalty} \\
        \midrule
        Data mapping           & Maintain records of processing    & 6 months  & Up to \$50,000 \\
        Privacy impact assess. & Conduct for high-risk processing  & 90 days   & Up to \$75,000 \\
        DPO appointment        & Designate for applicable orgs     & 6 months  & Up to \$25,000 \\
        Breach notification    & Report to NDPA within 72 hours    & Immediate & Up to \$200,000 \\
        Subject access rights  & Respond within 30 days            & Ongoing   & Up to \$10,000 \\
        Cross-border transfers & Complete transfer impact assess.  & 12 months & Up to \$100,000 \\
        \bottomrule
    \end{tabular}
    \caption{Summary of compliance obligations and penalties.}
\end{table}

\section{Enforcement}

\subsection{Investigation Powers}

The NDPA shall have the power to:

\begin{itemize}
    \item Conduct investigations and audits of any organization subject to
        this regulation.
    \item Issue binding orders requiring organizations to cease unlawful
        data processing.
    \item Order the restriction or suspension of data transfers.
    \item Compel the production of documents and records related to
        processing activities.
\end{itemize}

\subsection{Penalty Framework}

Violations of this regulation shall be penalized according to the
following graduated schedule:

\begin{description}
    \item[Level 1 --- Warning.] Issued for first-time minor violations
        with no evidence of harm. The organization shall receive a written
        warning and must demonstrate corrective action within 60 days.
    \item[Level 2 --- Administrative Fine.] Imposed for repeated or
        moderate violations. Fines range from \$10,000 to \$100,000
        depending on the nature and severity of the violation.
    \item[Level 3 --- Significant Fine.] Reserved for serious violations,
        including failure to report data breaches, systematic non-compliance,
        or processing without a lawful basis. Fines range from \$100,000
        to \$1,000,000.
    \item[Level 4 --- Maximum Sanction.] Imposed for egregious violations
        that cause substantial harm to data subjects. Fines may reach up
        to 4\% of the organization's annual global turnover or
        \$10,000,000, whichever is greater.
\end{description}

\section{Compliance Matrix}

\begin{table}[h]
    \centering
    \begin{tabular}{lcccc}
        \toprule
        \textbf{Article} & \textbf{Public Sector} & \textbf{Private (Large)} & \textbf{Processor} & \textbf{Cross-Border} \\
        \midrule
        Lawful basis         & \checkmark & \checkmark & ---     & \checkmark \\
        Data minimization    & \checkmark & \checkmark & \checkmark & \checkmark \\
        Storage limitation   & \checkmark & \checkmark & \checkmark & --- \\
        Privacy by design    & \checkmark & \checkmark & \checkmark & --- \\
        Breach notification  & \checkmark & \checkmark & \checkmark & \checkmark \\
        DPO appointment      & \checkmark & \checkmark & ---     & --- \\
        Subject access rights & \checkmark & \checkmark & ---     & \checkmark \\
        Transfer assessment  & \checkmark & \checkmark & \checkmark & \checkmark \\
        \bottomrule
    \end{tabular}
    \caption{Applicability of key articles by entity type.}
\end{table}

\section{Effective Date and Transitional Provisions}

This regulation shall enter into force on September 1, 2026. Organizations
shall have a transition period of 12 months from the effective date to
achieve full compliance. During the transition period, the NDPA shall
prioritize guidance and technical assistance over enforcement actions.

\printbibliography[heading=none]

\end{document}
